\textbf{Цели работы:}
\begin{itemize}
	\item познакомиться с инструментом Ansible;
	\item изучить формат описания конфигураций Ansible.
\end{itemize}

\textbf{Задачи:}
\begin{enumerate}
	\item изучить методические материалы к лабораторной работе;
	\item подготовить окружение;
	\item описать файл инвентаря с именем hosts, создав группы:
	\begin{itemize}
		\item staging\_serves, состоящий из двух хостов;
		\item prod\_serves, состоящий из одного хоста;
		\item all\_serves, объединяющий две предыдущие группы;
	\end{itemize}
	\item описать переменные для группы all\_serves;
	\item проверить доступность серверов через ad-hoc команду;
	\item написать Playbook для установки и запуска nginx.
\end{enumerate}

\section*{Выполнение задания}

\begin{figure}[!htbp]
	\centering
	\includegraphics[width=0.6\linewidth]{"Screenshot From 2026-04-20 00-57-49"}
	\caption{Установленный Ansible}
\end{figure}
\FloatBarrier

% cpu=1, ram=2048mb, rom>=5gb, adapter1.network_adapter=on, add gpt_partition=ext4='/'
% name=yeel, server_name=node-1, username='elanskaya-e-d'
% sudo apt update && sudo apt install opensh-server -y
% sudo systemctl enable ssh && sudo systemctl status ssh.service
% ip addr show
% повторить для node-2 и node-3
\begin{figure}[!htbp]
	\centering
	\includegraphics[width=0.3\linewidth]{"Screenshot (61)"}
	\caption{Созданные Ubuntu Server ВМ}
\end{figure}
\FloatBarrier
% yeel elanskayaed 1 % да, я не умею читать

\begin{figure}[!htbp]
	\centering
	\includegraphics[width=0.6\linewidth]{"Screenshot 2026-04-19 234031"}
	\caption{Параметры сервера}
\end{figure}
\FloatBarrier

\begin{figure}[!htbp]
	\centering
	\includegraphics[width=0.6\linewidth]{"Screenshot (62)"}
	\caption{Проверка доступности хостов}
\end{figure}
\FloatBarrier
% ansible all -m ping -i hosts -u elanskayaed --ask-pass
% sudo apt install sshpass

\begin{figure}[!htbp]
	\centering
	\includegraphics[width=0.6\linewidth]{"Screenshot (63)"}
	\caption{Запуск плейбука}
\end{figure}
\FloatBarrier
% ansible-playbook -i hosts nginx_install_and_run.yml -kK
% ansible-playbook nginx_install_and_run.yml -u elanskayaed --ask-pass

\begin{figure}[!htbp]
	\centering
	\includegraphics[width=0.6\linewidth]{"Screenshot (64)"}
	\caption{Приветственная страница nginx}
\end{figure}
\FloatBarrier
%\clearpage

\section*{Исходный код программы}

% ansible:
% hosts
% ansible.cfg
% nginx_install_and_run.yml
% group_vars/
% 	all_serves.yml
% roles/
% 	nginx/
% 		tasks/
% 			main.yml
% 		handlers/
% 			main.yml

\lstinputlisting[style=ansible, caption={Файл hosts}]{ansible/hosts}

\lstinputlisting[style=ansible, caption={Файл ansible.cfg}]{ansible/ansible.cfg}

\lstinputlisting[style=ansible, caption={Файл для переменных группы all\_serves}]{ansible/group_vars/all_serves.yml}

\lstinputlisting[style=ansible, caption={Плейбук nginx\_install\_and\_run.yml}]{ansible/nginx_install_and_run.yml}

\lstinputlisting[style=ansible, caption={Файл задачи для роли nginx}]{ansible/roles/nginx/tasks/main.yml}

\lstinputlisting[style=ansible, caption={Файл обработчика для роли nginx}]{ansible/roles/nginx/handlers/main.yml}

\section*{Заключение}

В ходе выполнения лабораторной работы были решены следующие задачи:
\begin{itemize}
	\item проведено ознакомление с инструментом Ansible;
	\item изучен формат описания конфигураций Ansible.
\end{itemize}

\section*{Контрольные вопросы}

\begin{enumerate}
	\item \textit{Что такое Ansible?} Система управления конфигурациями.
	\item \textit{Что такое Playbook?} Повторяющаяся, многократно используемая, простая система управления конфигурацией и развертывания на нескольких машинах, которая хорошо подходит для развертывания сложных приложений.
	\item \textit{Для чего нужны роли?} Роли позволяют автоматически загружать связанные переменные, файлы, задачи, обработчики и другие артефакты Ansible на основе известной структуры файлов. Сгруппировав содержимое по ролям, можно использовать их повторно и делиться ими с другими пользователями.
	\item \textit{Что гарантирует идемпотентность в Ansible?} Плейбуки и роли.
	% Идемпотентность -- свойство объекта или операции при повторном применении операции с одними и теми же параметрами давать один и тот же результат.
\end{enumerate}
